% GENERATED FILE. Edit canonical lessons, then run npm run generate:books.
% canonical-source-hash: fnv1a64:d1d99bdd3842b434
% canonical-lessons: TE-C398-cattam, TE-C398-ullanghincu, TE-C398-amalu-ceyu, TE-C398-siksa, TE-C398-nisedhincu

\chapter{Laws and Rules}
\label{ch:te-laws-and-rules}
\hlchaptermodality{398}

\begin{chapteropening}
\textbf{By the end of this chapter:} \emph{I can talk about laws, breaking them and punishment.}

Complete the last lesson of chapter 398: I can talk about laws, breaking them and punishment.
\end{chapteropening}

\section[\texorpdfstring{\te{చట్టం} (caṭṭaṁ)}{చట్టం (caṭṭaṁ)}]{\te{చట్టం} (caṭṭaṁ) — a law; an act}

\label{lesson:TE-C398-cattam}



\begin{quote}
Before the new one: say the Telugu for to mention, then the Telugu for to express.
\end{quote}

\subsection*{You'll want to know: \te{చట్టం}}
\textbf{\te{చట్టం}} — \emph{caṭṭaṁ} — \textquotedblleft{}a law; an act\textquotedblright{}.

\begin{quote}
\te{ఆ} \te{పని} \te{చట్టం} \te{ప్రకారం} \te{తప్పు}.
\end{quote}

\begin{quote}
\emph{ā pani caṭṭaṁ prakāraṁ tappu.}
\end{quote}

\begin{quote}
That act is wrong under the law.
\end{quote}

Plural: \te{చట్టాలు}, \emph{caṭṭālu}.

\te{చట్టం} \te{ప్రకారం}, according to the law.

\begin{cousinweb}
\textbf{\te{చట్టం}} (\emph{caṭṭaṁ}) is native Dravidian: DEDR 2304 lists Telugu \emph{caṭṭamu}, \textquotedblleft{}a frame, a plan, a system, a law\textquotedblright{}, with Tamil \emph{caṭṭam}, \textquotedblleft{}a frame; a rule, a regulation\textquotedblright{}.
\end{cousinweb}

\begin{grammarlens}[title={What you've built}]
One more word for giving reasons and opinions.
\end{grammarlens}

\subsection*{Your turn}
\begin{itemize}
\raggedright
  \item \emph{Say it:} \emph{caṭṭaṁ}
  \item \emph{Say it:} \emph{caṭṭaṁ}, once more
  \item \emph{Say it:} \emph{ā pani caṭṭaṁ prakāraṁ tappu.} — That act is wrong under the law.
  \item \emph{Recall:} say the Telugu for to mention, then the Telugu for to express, then say \emph{caṭṭaṁ} again
\end{itemize}

\begin{tcolorbox}[breakable,colback=teal!4,colframe=teal!35!black,title={Before you move on}]
What does \te{చట్టం} mean? (\textquotedblleft{}A law; an act\textquotedblright{}.) Say it once more.
\end{tcolorbox}
\section[\texorpdfstring{\te{ఉల్లంఘించు} (ullaṅghiñcu)}{ఉల్లంఘించు (ullaṅghiñcu)}]{\te{ఉల్లంఘించు} (ullaṅghiñcu) — to violate; to break (a rule)}

\label{lesson:TE-C398-ullanghincu}



\begin{quote}
Before the new one: say the Telugu for to express, then the Telugu for a law.
\end{quote}

\subsection*{You'll want to know: \te{ఉల్లంఘించు}}
\textbf{\te{ఉల్లంఘించు}} — \emph{ullaṅghiñcu} — \textquotedblleft{}to violate; to break (a rule)\textquotedblright{}.

\begin{quote}
\te{అతను} \te{కంపెనీ} \te{నియమాలను} \te{ఉల్లంఘించాడు}.
\end{quote}

\begin{quote}
\emph{atanu kampenī niyamālanu ullaṅghiñcāḍu.}
\end{quote}

\begin{quote}
He broke the company rules.
\end{quote}

Past (I …): \te{ఉల్లంఘించాను}, \emph{ullaṅghiñcānu}.

Formal: for rules, laws and agreements.

\begin{cousinweb}
\textbf{\te{ఉల్లంఘించు}} (\emph{ullaṅghiñcu}) is built on Sanskrit \emph{ullaṅghana}, \textquotedblleft{}the act of leaping over\textquotedblright{} (CDIAL 2365), and so \textquotedblleft{}transgression\textquotedblright{} (Monier-Williams), with the Telugu verb ending -\te{ఇంచు}.
\end{cousinweb}

\begin{grammarlens}[title={What you've built}]
One more word for giving reasons and opinions.
\end{grammarlens}

\subsection*{Your turn}
\begin{itemize}
\raggedright
  \item \emph{Say it:} \emph{ullaṅghiñcu}
  \item \emph{Say it:} \emph{ullaṅghiñcu}, once more
  \item \emph{Say it:} \emph{atanu kampenī niyamālanu ullaṅghiñcāḍu.} — He broke the company rules.
  \item \emph{Recall:} say the Telugu for to express, then the Telugu for a law, then say \emph{ullaṅghiñcu} again
\end{itemize}

\begin{tcolorbox}[breakable,colback=teal!4,colframe=teal!35!black,title={Before you move on}]
What does \te{ఉల్లంఘించు} mean? (\textquotedblleft{}To violate; to break (a rule)\textquotedblright{}.) Say it once more.
\end{tcolorbox}
\section[\texorpdfstring{\te{అమలు} \te{చేయు} (amalu cēyu)}{అమలు చేయు (amalu cēyu)}]{\te{అమలు} \te{చేయు} (amalu cēyu) — to implement; to enforce}

\label{lesson:TE-C398-amalu-ceyu}



\begin{quote}
Before the new one: say the Telugu for a law, then the Telugu for to violate.
\end{quote}

\subsection*{You'll want to know: \te{అమలు} \te{చేయు}}
\textbf{\te{అమలు} \te{చేయు}} — \emph{amalu cēyu} — \textquotedblleft{}to implement; to enforce\textquotedblright{}.

\begin{quote}
\te{ప్రభుత్వం} \te{ఆ} \te{చట్టాన్ని} \te{వెంటనే} \te{అమలు} \te{చేసింది}.
\end{quote}

\begin{quote}
\emph{prabhutvaṁ ā caṭṭānni veṇṭanē amalu cēsindi.}
\end{quote}

\begin{quote}
The government put that law into effect at once.
\end{quote}

Past (I …): \te{అమలు} \te{చేశాను}, \emph{amalu cēśānu}.

\te{అమలు} alone is enforcement: \te{అమలులోకి} \te{వచ్చు}, to come into force.

\begin{cousinweb}
\textbf{\te{అమలు} \te{చేయు}} (\emph{amalu cēyu}) is usually traced to Hindustani \emph{amal}, \textquotedblleft{}action, execution\textquotedblright{}, a word of Arabic origin, with the native verb \te{చేయు}, \textquotedblleft{}to do\textquotedblright{} (DEDR 1957).
\end{cousinweb}

\begin{grammarlens}[title={What you've built}]
One more word for giving reasons and opinions.
\end{grammarlens}

\subsection*{Your turn}
\begin{itemize}
\raggedright
  \item \emph{Say it:} \emph{amalu cēyu}
  \item \emph{Say it:} \emph{amalu cēyu}, once more
  \item \emph{Say it:} \emph{prabhutvaṁ ā caṭṭānni veṇṭanē amalu cēsindi.} — The government put that law into effect at once.
  \item \emph{Recall:} say the Telugu for a law, then the Telugu for to violate, then say \emph{amalu cēyu} again
\end{itemize}

\begin{tcolorbox}[breakable,colback=teal!4,colframe=teal!35!black,title={Before you move on}]
What does \te{అమలు} \te{చేయు} mean? (\textquotedblleft{}To implement; to enforce\textquotedblright{}.) Say it once more.
\end{tcolorbox}
\section[\texorpdfstring{\te{శిక్ష} (śikṣa)}{శిక్ష (śikṣa)}]{\te{శిక్ష} (śikṣa) — a punishment; a sentence}

\label{lesson:TE-C398-siksa}



\begin{quote}
Before the new one: say the Telugu for to violate, then the Telugu for to implement.
\end{quote}

\subsection*{You'll want to know: \te{శిక్ష}}
\textbf{\te{శిక్ష}} — \emph{śikṣa} — \textquotedblleft{}a punishment; a sentence\textquotedblright{}.

\begin{quote}
\te{దొంగకు} \te{రెండు} \te{సంవత్సరాల} \te{శిక్ష} \te{పడింది}.
\end{quote}

\begin{quote}
\emph{doṅgaku reṇḍu saṁvatsarāla śikṣa paḍindi.}
\end{quote}

\begin{quote}
The thief got a two-year sentence.
\end{quote}

Plural: \te{శిక్షలు}, \emph{śikṣalu}.

\te{శిక్ష} \te{పడు}, to be punished (a punishment falls on someone); \te{శిక్షించు}, to punish.

\begin{cousinweb}
\textbf{\te{శిక్ష}} (\emph{śikṣa}) is borrowed from Sanskrit \emph{śikṣā}, \textquotedblleft{}learning, study\textquotedblright{}, and later also \textquotedblleft{}punishment\textquotedblright{} (CDIAL 12432). Telugu shortens the final long a.
\end{cousinweb}

\begin{grammarlens}[title={What you've built}]
One more word for giving reasons and opinions.
\end{grammarlens}

\subsection*{Your turn}
\begin{itemize}
\raggedright
  \item \emph{Say it:} \emph{śikṣa}
  \item \emph{Say it:} \emph{śikṣa}, once more
  \item \emph{Say it:} \emph{doṅgaku reṇḍu saṁvatsarāla śikṣa paḍindi.} — The thief got a two-year sentence.
  \item \emph{Recall:} say the Telugu for to violate, then the Telugu for to implement, then say \emph{śikṣa} again
\end{itemize}

\begin{tcolorbox}[breakable,colback=teal!4,colframe=teal!35!black,title={Before you move on}]
What does \te{శిక్ష} mean? (\textquotedblleft{}A punishment; a sentence\textquotedblright{}.) Say it once more.
\end{tcolorbox}
\section[\texorpdfstring{\te{నిషేధించు} (niṣēdhiñcu)}{నిషేధించు (niṣēdhiñcu)}]{\te{నిషేధించు} (niṣēdhiñcu) — to ban; to prohibit}

\label{lesson:TE-C398-nisedhincu}



\begin{quote}
Before the new one: say the Telugu for to implement, then the Telugu for a punishment.
\end{quote}

\subsection*{You'll want to know: \te{నిషేధించు}}
\textbf{\te{నిషేధించు}} — \emph{niṣēdhiñcu} — \textquotedblleft{}to ban; to prohibit\textquotedblright{}.

\begin{quote}
\te{పాఠశాలలో} \te{ఫోన్లను} \te{నిషేధించారు}.
\end{quote}

\begin{quote}
\emph{pāṭhaśālalō phōnlanu niṣēdhiñcāru.}
\end{quote}

\begin{quote}
They banned phones at the school.
\end{quote}

Past (I …): \te{నిషేధించాను}, \emph{niṣēdhiñcānu}.

The noun is \te{నిషేధం}, a ban.

\begin{cousinweb}
\textbf{\te{నిషేధించు}} (\emph{niṣēdhiñcu}) is built on Sanskrit \emph{niṣedha}, \textquotedblleft{}warding off, hindering, prohibition\textquotedblright{} (Monier-Williams), with the Telugu verb ending -\te{ఇంచు}.
\end{cousinweb}

\begin{grammarlens}[title={What you've built}]
One more word for giving reasons and opinions.
\end{grammarlens}

\subsection*{Your turn}
\begin{itemize}
\raggedright
  \item \emph{Say it:} \emph{niṣēdhiñcu}
  \item \emph{Say it:} \emph{niṣēdhiñcu}, once more
  \item \emph{Say it:} \emph{pāṭhaśālalō phōnlanu niṣēdhiñcāru.} — They banned phones at the school.
  \item \emph{Recall:} say the Telugu for to implement, then the Telugu for a punishment, then say \emph{niṣēdhiñcu} again
\end{itemize}

\begin{tcolorbox}[breakable,colback=teal!4,colframe=teal!35!black,title={Before you move on}]
What does \te{నిషేధించు} mean? (\textquotedblleft{}To ban; to prohibit\textquotedblright{}.) Say it once more.
\end{tcolorbox}
